% Generated draft: requires mathematical and source review.
\documentclass[11pt]{article}
\usepackage[margin=25mm]{geometry}
\usepackage{fontspec}
\setmainfont{Latin Modern Roman}
\usepackage{amsmath,amssymb,mathtools,unicode-math}
\setmathfont{Latin Modern Math}
\usepackage{enumitem,needspace,xurl,hyperref}
\setlength{\parindent}{0pt}
\setlength{\parskip}{5pt}
\setlength{\emergencystretch}{3em}
\newcommand{\N}{\mathbb{N}}
\newcommand{\Z}{\mathbb{Z}}
\newcommand{\Q}{\mathbb{Q}}
\newcommand{\R}{\mathbb{R}}
\newcommand{\C}{\mathbb{C}}
\newcommand{\F}{\mathbb{F}}
\newcommand{\A}{\mathbb{A}}
\newcommand{\PP}{\mathbb{P}}
\newcommand{\E}{\mathbb{E}}
\newcommand{\eps}{\varepsilon}
\newcommand{\dd}{\,\mathrm{d}}
\begin{document}
\section*{UnsolvedMath ID: 30004505}
\textbf{Source problem code:} catalogue entry
\subsection{Definitions and mathematical statement}
The supplied formulation requires conventions for the degree of a nilmanifold and for polynomial sequences. The following gives a precise filtered-nilmanifold interpretation; the record does not specify whether this is its intended convention.

Let \(k\geq 1\) be an integer. Let \(G\) be a nilpotent Lie group and \(\Gamma\leq G\) a discrete cocompact subgroup, so that \(G/\Gamma\) is compact. Fix a metric \(d\) on \(G/\Gamma\). Equip \(G\) with a descending filtration by closed subgroups
\[
G=G_0=G_1\supseteq G_2\supseteq\cdots\supseteq G_k\supseteq G_{k+1}=\{e\},
\qquad [G_i,G_j]\subseteq G_{i+j},
\]
where \(G_r=\{e\}\) for \(r>k\). This is the degree-at-most-\(k\) convention used here; it is additional interpretive data absent from the record. For \(g:\Z\to G\), define
\[
D_hg(n)=g(n+h)g(n)^{-1}.
\]
Write \(\operatorname{Poly}(\Z\to G_\bullet)\) for the sequences satisfying
\[
D_{h_r}\cdots D_{h_1}g(n)\in G_r
\quad\text{for every }r\geq1\text{ and all }n,h_1,\ldots,h_r\in\Z.
\]
The source writes \(\operatorname{Poly}(\Z\to G)\) without specifying a filtration.

For a positive integer \(n\), the M\"obius function \(\mu(n)\) is \(0\) if a prime square divides \(n\), and otherwise is \((-1)^r\) if \(n\) has \(r\) distinct prime factors, with \(\mu(1)=1\). A map \(F:G/\Gamma\to\C\) is Lipschitz if \(|F(u)-F(v)|\leq Ld(u,v)\) for some finite \(L\) and all \(u,v\).

Under the preceding convention, the proposed assertion is as follows. For every fixed such filtered nilmanifold and every fixed Lipschitz \(F\), and for every positive function \(H=H(X)\) with \(H(X)\to\infty\) as \(X\to\infty\), define
\[
A_{X,H}(x)=\sup_{g\in\operatorname{Poly}(\Z\to G_\bullet)}
\left|\sum_{\substack{n\in\Z_{\geq1}\\x\leq n\leq x+H(X)}}\mu(n)F(g(n)\Gamma)\right|.
\]
Is it true that, for every \(\varepsilon>0\),
\[
\lim_{X\to\infty}\frac{1}{X}\lambda\bigl(\{x\in[X,2X]:A_{X,H}(x)>\varepsilon H(X)\}\bigr)=0,
\]
where \(\lambda\) denotes Lebesgue measure? This specifies one meaning of the source's assertion \(A_{X,H}(x)=o(H)\) for almost all \(x\in[X,2X]\). If integer starting points are intended, replace the measure in this display by the cardinality of the corresponding set with \(x\in\Z\). The record does not choose between these interpretations. Both endpoints of the sum are included, as in the supplied statement. No upper bound on \(H(X)\) is supplied, and none is imposed here. The supremum allows the polynomial sequence to vary with \(X\) and \(x\), while the nilmanifold, filtration, metric, and \(F\) remain fixed.
\subsection{Short English statement}
The question asks whether M\"obius sums on intervals of every diverging length have negligible correlation, relative to the interval length, uniformly over polynomial nilsequences on a fixed nilmanifold, outside a vanishing proportion of starting points. The filtration convention and the meaning of almost all starting points require confirmation from the source.
\subsection{Sources}
\begin{itemize}
\item Supplied local UnsolvedMath registry record, permanent ID \(30004505\), title ``Short-Interval M\"obius Orthogonality to Nilsequences,'' source problem code \texttt{catalogue entry}, in the Oberwolfach Reports Open Problems collection. The record gives proposed year 2020 and names no proposer.
\item Supplied source citation: \emph{Automorphic Forms and Arithmetic} (2021), Oberwolfach Reports, pp.\ 1317--1383, DOI \texttt{10.4171/owr/2020/26}. The literature-source label is ``Oberwolfach Report 26/2020.'' \url{https://doi.org/10.4171/owr/2020/26}. The supplied background describes these reports as documenting mathematical workshops and their open-problem discussions.
\item Supplied literature source: \emph{Higher uniformity of bounded multiplicative functions}. \url{https://arxiv.org/abs/2007.15644}.
\item Supplied literature source: \emph{Higher uniformity of arithmetic functions II}. \url{https://arxiv.org/abs/2411.05770}.
\item Literature assessment dated exactly \texttt{2026-08-22}: the supplied curation labels the problem open and states that the full nilsequence uniformity conjecture for arbitrary diverging interval lengths remains open, while a main theorem establishes a substantial power-scale range but not the stated endpoint. This is supplied dated curation, not an independently verified literature assessment. The background explicitly requests independent verification of time-sensitive status and cited claims.
\item The supplied statement-status label is \texttt{corrected\_verified}; its accompanying explanation says that mathematical content was preserved while notation, typography, omitted definitions, or duplicate clauses were repaired. This is a registry qualification, not independent mathematical verification of this draft.
\end{itemize}
\end{document}
